\begin{lemma}[The differentiated equation]
Let $1<p<2$, let $U\subseteq\mathrm{Plane}$ and let $v\colon\mathrm{Plane}\to\mathbb{R}$.
Assume:
\begin{enumerate}
\item[(H1)] $v$ is $p$-harmonic on $U$ with weak gradient its own classical gradient $Dv$, in
the sense of \noderef{IsPHarm}; in particular $U$ is open;
\item[(H2)] $v$ is of class $C^{\infty}$ on $U$;
\item[(H3)] $Dv(x)\ne0$ for every $x\in U$.
\end{enumerate}
Let $e\in\mathrm{Plane}$ and write $\partial_ev(x):=\langle Dv(x),e\rangle$ for the directional
derivative of $v$ in the direction $e$.  Then, in the sense of \noderef{IsLinearDivSol},
\[
  \partial_ev\ \text{ is a weak solution of }\
  \operatorname{div}\bigl(A\,D(\partial_ev)\bigr)=0\ \text{ in }U\
  \text{ with weak gradient }D(\partial_ev),
  \qquad A(x):=D_\xi A_p\bigl(Dv(x)\bigr),
\]
where $A_p$ is the $p$-flux of \noderef{PFlux} and $D_\xi A_p$ the flux Jacobian of
\noderef{FluxJacobian}.  Explicitly, the five conditions of \noderef{IsLinearDivSol} hold for
the quadruple $\bigl(U,\;x\mapsto D_\xi A_p(Dv(x)),\;\partial_ev,\;D(\partial_ev)\bigr)$.

Taking $e$ to be the two vectors of the standard basis recovers the paper's two differentiated
equations $\operatorname{div}\bigl(D_\xi A_p(Dv)\,D(v_{x_i})\bigr)=0$ in $B_{2\tau}$, $i=1,2$
(paper line 177).  No pointwise divergence operator appears: the assertion is the weak
formulation of that identity.
\end{lemma}
\begin{proof}
Throughout, $U$ is open by (H1).  Write
\[
  \xi(x):=Dv(x),\qquad w:=\partial_ev,\qquad G:=Dw,\qquad A(x):=D_\xi A_p\bigl(\xi(x)\bigr).
\]

We will use \noderef{FluxJacobianElliptic} only through its clauses (1) and (6), which are
statements about all $\xi\ne0$ and about the whole open set $\{\zeta\ne0\}$ and so do not
involve its annulus parameters; we invoke it with $1<p<2$ and $a=b=1$, which satisfy
$0<a$ and $a\le b$.  Those clauses give:
\begin{itemize}
\item[(E1)] for every $\zeta\ne0$, $A_p$ is differentiable at $\zeta$ with
  $DA_p(\zeta)=D_\xi A_p(\zeta)$;
\item[(E2)] $A_p$ is of class $C^{\infty}$ on the open set $\{\zeta\ne0\}$ (clause (6)).
\end{itemize}

\medskip\noindent\textbf{Step 1: regularity.}
By (H2) $v$ is of class $C^{n+1}$ on $U$ for every $n$.  Hence, by clause (1) of the background
proposition on regularity of the gradient field and of directional derivatives, the gradient field
$\xi=Dv\colon\mathrm{Plane}\to\mathrm{Plane}$ is of class $C^{\infty}$ on $U$; in particular it
is continuous on $U$ and, by the meaning of $C^{1}$ regularity, differentiable at every point
of $U$.  By clause (2) of the same
proposition, $w=\partial_ev$ is of class $C^{\infty}$ on $U$; in particular $w$ is continuous on
$U$ and, by clause (1) applied to $w$, $G=Dw$ is of class $C^{\infty}$ on $U$, hence continuous
on $U$.

By (H3), $\xi(U)\subseteq\{\zeta\ne0\}$, which is open.  By (E2) and clause (3) of the
background proposition on the chain rule for compositions (applied with $H=\xi$ on $U$ and
$\Phi=A_p$ on $\{\zeta\ne0\}$), the composite
\[
  X:=A_p\circ\xi
\]
is of class $C^{\infty}$ on $U$; in particular $X$ is continuous on $U$, is differentiable at
every point of $U$ (by the meaning of $C^{1}$ regularity), and $x\mapsto DX(x)$ is continuous
on $U$ by clause (5) of the background proposition on the recursion for the iterated Fr\'echet
derivative, applied with $n=1$ and $W=\mathrm{Plane}$.  Consequently
$\partial_eX(x)=DX(x)[e]$ is continuous on $U$, evaluation at $e$ being a continuous linear map
on the space of continuous linear forms.

\medskip\noindent\textbf{Step 2: the chain rule in the direction $e$.}
Fix $x\in U$.  By Step 1, $\xi$ is differentiable at $x$; by (E1) and (H3), $A_p$ is
differentiable at $\xi(x)$ with derivative $D_\xi A_p(\xi(x))=A(x)$.  Clauses (1) and (2) of
the chain-rule proposition therefore give
\begin{equation}\label{eq:de-chain}
  \partial_eX(x)=DX(x)[e]=A(x)\bigl[D\xi(x)[e]\bigr]=A(x)\bigl[\partial_e(Dv)(x)\bigr].
\end{equation}
Since $v$ is of class $C^{2}$ on $U$, clause (3) of the gradient-field proposition
gives $\partial_e(Dv)(x)=D(\partial_ev)(x)=G(x)$.  Substituting into \eqref{eq:de-chain},
\begin{equation}\label{eq:de-key}
  \partial_eX(x)=A(x)\bigl[G(x)\bigr]\qquad\text{for every }x\in U.
\end{equation}

\medskip\noindent\textbf{Step 3: conditions (1), (2) and (3) of \noderef{IsLinearDivSol}.}

\emph{(1) Continuity.}  $w$ is continuous on $U$ by Step 1.

\emph{(2) Measurability.}  $U$ is Lebesgue measurable (clause (2) of the integrability
proposition) and $G$ is continuous on $U$ by Step 1, so $G$ is almost everywhere strongly
measurable with respect to Lebesgue measure restricted to $U$, by clause (3) of that
proposition.

\emph{(3) Local $L^{2}$ integrability.}  The function $x\mapsto\lvert G(x)\rvert^{2}$ is
continuous on $U$, being the composition of the continuous $G$ with the continuous map
$z\mapsto\lvert z\rvert^{2}$.  By clause (3) of the integrability proposition it is therefore
bounded, hence integrable, on every compact subset of $U$.

\medskip\noindent\textbf{Step 4: condition (4), $G$ is the distributional gradient of $w$.}
This is \noderef{ClassicalGradIsWeakGrad}, applied with the open set $U$ and the function $w$.
Its two hypotheses hold: $U$ is open by (H1), and $w$ is of class $C^{1}$ on $U$ by Step 1,
which gives $w$ of class $C^{\infty}$ there.  Its conclusion is that for every
$\varphi\colon\mathrm{Plane}\to\mathbb{R}$ of class $C^{\infty}$ with compact support and
$\operatorname{tsupp}\varphi\subseteq U$,
\[
  \int_U w\,D\varphi\,dx=-\int_U\varphi\,Dw\,dx=-\int_U\varphi\,G\,dx ,
\]
the last equality because $G=Dw$ by definition.  That is condition (4).

\medskip\noindent\textbf{Step 5: condition (5), the weak equation.}
Let $\psi\colon\mathrm{Plane}\to\mathbb{R}$ be of class $C^{\infty}$ with compact support and
$\operatorname{tsupp}\psi\subseteq U$.  We must show
\begin{equation}\label{eq:de-goal}
  \int_U\bigl\langle A(x)\,G(x),\,D\psi(x)\bigr\rangle\,dx=0 .
\end{equation}

\emph{(5a) $\partial_e\psi$ is again a test function.}  Put $\varphi:=\partial_e\psi$.  By
clause (4) of the gradient-field proposition, $\varphi$ is of class $C^{\infty}$ on
$\mathrm{Plane}$, has compact support, and
$\operatorname{tsupp}\varphi\subseteq\operatorname{tsupp}\psi\subseteq U$.  Hypothesis (H1) and
the weak equation of \noderef{IsPHarm}, applied to the test function $\varphi$, therefore give
\begin{equation}\label{eq:de-weak}
  \int_U\bigl\langle A_p\bigl(Dv(x)\bigr),\,D\varphi(x)\bigr\rangle\,dx=0,
  \qquad\text{i.e.}\qquad
  \int_U\bigl\langle X(x),\,D\varphi(x)\bigr\rangle\,dx=0 .
\end{equation}
Moreover, since $\psi$ is of class $C^{2}$ on the open set $\mathrm{Plane}$, clause (3) of the
gradient-field proposition gives
\begin{equation}\label{eq:de-sym}
  D\varphi(x)=D(\partial_e\psi)(x)=\partial_e\bigl(D\psi\bigr)(x)
  \qquad(x\in\mathrm{Plane}).
\end{equation}

\emph{(5b) $D\psi$ is admissible in the integration-by-parts lemma.}  Put $Y:=D\psi$.  By
clause (1) of the gradient-field proposition, applied with $\Omega=\mathrm{Plane}$ and $f=\psi$
of class $C^{\infty}$ there, $Y$ is of class $C^{\infty}$ --- hence of class $C^{1}$ --- on all
of $\mathrm{Plane}$.  By clause (3) of the background proposition on the product rule and
linearity for the classical gradient, $Y=D\psi$ vanishes on the open set
$\mathrm{Plane}\setminus\operatorname{tsupp}\psi$, on which $\psi\equiv0$; hence the support of
$Y$ is contained in $\operatorname{tsupp}\psi$, and that set being closed,
$\operatorname{tsupp}Y\subseteq\operatorname{tsupp}\psi\subseteq U$.  In particular
$\operatorname{tsupp}Y$ is a closed subset of the compact set $\operatorname{tsupp}\psi$, hence
compact, so $Y$ has compact support.

\emph{(5c) Integration by parts in the direction $e$.}  We may now apply
\noderef{DirectionalIntegrationByParts} with the open set $U$, the field $X=A_p\circ\xi$, which
is of class $C^{1}$ on $U$ by Step 1, the field $Y=D\psi$, which by (5b) is of class $C^{1}$ on
all of $\mathrm{Plane}$ with compact support and $\operatorname{tsupp}Y\subseteq U$, and the
direction $e$.  It gives
\begin{equation}\label{eq:de-ibp}
  \int_U\bigl\langle\partial_eX(x),\,D\psi(x)\bigr\rangle\,dx
  =-\int_U\bigl\langle X(x),\,\partial_e\bigl(D\psi\bigr)(x)\bigr\rangle\,dx .
\end{equation}
By \eqref{eq:de-sym} the integrand on the right-hand side equals
$\langle X(x),D\varphi(x)\rangle$ at every $x\in\mathrm{Plane}$, so the two integrals over $U$
agree by clause (5) of the background proposition on restriction and congruence; and that
integral vanishes by \eqref{eq:de-weak}.  The right-hand side of \eqref{eq:de-ibp} is therefore
$0$, and so
\begin{equation}\label{eq:de-half}
  \int_U\bigl\langle\partial_eX(x),\,D\psi(x)\bigr\rangle\,dx=0 .
\end{equation}

\emph{(5d) Conclusion.}  By \eqref{eq:de-key} the integrands
$\langle\partial_eX(x),D\psi(x)\rangle$ and $\langle A(x)G(x),D\psi(x)\rangle$ are equal at
every $x\in U$, so their integrals over $U$ agree, again by clause (5) of the
restriction-and-congruence proposition.  With \eqref{eq:de-half} this is \eqref{eq:de-goal},
i.e.\ condition (5), and the proof is complete.
\end{proof}
